\clearpage
\part*{Part XIV. Condensed Matter And Atoms}
\label{atlaspart:14-condensed-matter-and-atoms}
\addcontentsline{toc}{part}{Part XIV. Condensed Matter And Atoms}
\section*{Introduction to Part XIV}\label{narr:part:XII}
\addcontentsline{toc}{section}{Introduction to Part XIV}
A lattice is a periodic sink. Part XIV reads condensed matter and atomic physics from that one reading: the chain of copies and the periodic readout, the band at the periodic web, the chain of pairs and the charged periodic sink, the winding of the wash holonomy that the literature calls the Chern number, the periodic table as the joined-electron sinks, and the atom as the hydrogen probe web with its many-body readout. The Part draws on the probe and the saturation of Part I, on the measurement of Part III, on the hadron chapters of Part IX and on the lepton ladder of Part VII, and it is one of the Parts with the most theorems, since a periodic web is an object the saturation describes exactly.

The order of the six chapters runs from the smallest periodic object to the largest. A ring of copies of one body joined at their docks is the simplest periodic web, and chapter 12-1 builds it, asks which rings the grammar admits and finds the sink the pairing move reaches. Chapter 12-2 reads the spectrum of the step generator on such a ring, which is what a band is here, and chapter 12-3 reads the chain of pairs that the join links of a ring form. Chapter 12-4 reads a winding, the holonomy of the wash carried once around a period and around the full cycle of a joined network, which is where the integers of the quantum Hall literature appear as counts. The last two chapters join many electrons to a probe and many bodies to their measuring bodies, and both are open at the box. A reader who wants to see why a lattice has a band without a Bloch theorem should read chapters 12-1 and 12-2 together: the translation is an automorphism of the ring, the step generator is tiled by it, and the spectrum follows from the cycle order.

Chapter 12-1 builds the lattice. Theorem~224.1 reads the chain of copies and the periodic readout, Theorem~224.2 the periodic network admitted by the grammar, and Theorem~224.3 the periodic sink at a named cell count, three theorems, with the periodic network of hub assemblies, Open result~224.4, open beside them. A chain of two, three or four copies of the electron body joined at the dock by the pairing move, the join graph a path or a cycle, is one network whose direct record sum is an exact element at every one of the five arrays, every record graded five ways, and the readout is periodic under the cycle shift, Theorem~224.1. The periodic network is admitted by the grammar at both cells of the item, the library star with its dock, where rings of docked stars at three dock kinds have every join link a (2,2) pairing link between two docks and the translation an automorphism, fifty-four networks built, seventeen periodic and twenty-nine admitted, and the two-dock qutrit, Theorem~224.2. The periodic sink exists at a named cell count with its cell an output: at the library star the sink of the ring of docked stars under the pairing move is periodic at the even counts two, four and six and refused at the odd counts three and five, where the charge is carried at one open dock, Theorem~224.3. What stands open is the periodic network of hub assemblies, the first link refused at two cells being the (3,3) join link between two hub assemblies, neither end a dock qubit, Open result~224.4. Chapter 12-2 reads the band at the periodic web in Theorem~225.1, a theorem resting on the records, the probe, the energy chapter and the lattice: at the rings of docked stars and of docked qutrits re-built by the same join at two to six cells, every link is a register of the step generator tiled by the translation, and the cycle-order spectrum of the ring is the spectrum the literature calls the band, Theorem~225.1.

Chapter 12-3 reads the chain of pairs at the periodic web, Theorem~226.1, a theorem, and states two open results together in Open result~226.2, the pairing of two charged bodies and the charged periodic sink. At every ring of the lattice, the twelve rings of the lattice tables and sixteen more built at five to eight cells, the join links are (2,2) pairing links between dock-labeled pendant qubits of consecutive copies, disjoint and carried to the next by the translation, so the chain of pairs is admitted at every ring, Theorem~226.1. The pairing of two charged bodies at the dock qubits is refused at the two-dock qutrit at two cells, the link between the electron dock and the first dock of the first cell lying beyond the valence bound with every dock of every cell consumed by the ring join, which is the first parting, Open result~226.2. Chapter 12-4 reads the winding: the table titled The Chern number, awaited (in chapter~\ref{mod:XII-4-chern-number}, among the instruments) holds the chapter as awaited, the theorem titled The winding of the wash holonomy around one period (in chapter~\ref{mod:XII-4-chern-number}, among the instruments) reads the winding of the wash holonomy around one period and the theorem titled The winding around the full cycle of a joined periodic network (in chapter~\ref{mod:XII-4-chern-number}, among the instruments) around the full cycle of a joined periodic network, both theorems, the open result titled The flux quantization premise re-posed (in chapter~\ref{mod:XII-4-chern-number}, among the instruments) re-poses the flux quantization premise as an open result, and the theorem titled The flux clause re-posed at the face (chapter~\ref{mod:XII-4-D538575-flxf-t1}, among the instruments) re-poses the flux clause at the face, a theorem. At every periodic stable configuration the winding of the wash holonomy around one period, the holonomy a map on the register letters, is an exact integer, twelve twenty-fourths of a turn at every one of the twenty-eight deciding docks, the theorem titled The winding of the wash holonomy around one period (in chapter~\ref{mod:XII-4-chern-number}, among the instruments). The winding around the full cycle of every joined periodic network, lifted per loop, equals half the count of its parity-odd stable configurations at 175 of 175 networks, the nonzero integer reached at 133 networks beyond one sink, with the quantum Hall integer a name beside, the theorem titled The winding around the full cycle of a joined periodic network (in chapter~\ref{mod:XII-4-chern-number}, among the instruments). The flux integer per face is the fan sum of that face at every two-path network, at every one of the 45739 defined two-path networks read as faces with the enclosed set decomposed into its 89607 fan loops, the theorem titled The flux clause re-posed at the face (chapter~\ref{mod:XII-4-D538575-flxf-t1}, among the instruments). The reader who knows the Chern number as an integral of a curvature over a torus will find it here as a count of parity-odd configurations, and the box is explicit that the quantum Hall integer is a name beside the count.

Chapters 12-5 and 12-6 read atoms. Open result~228.1 states the periodic table as the joined-electron sinks as an open result, resting on the moves, the saturation, the probe, the corollaries of Part III and the hydrogen chapter, and Theorem~229.1 states the probe read at the four bodies joined to the measuring body, likewise open. At the reading taken as the one of reference, the probe at its three docks with the pairable set the docks and the saturation across copies, the join of library electrons at the probe is admitted at every count from one to ten and refused beyond, which is where the box names its parting, Open result~228.1. The four bodies joined to their measuring bodies at the attachment site, with the joined read opened at the dock or nowhere under the transport of reference, give the unit and the register as outputs of the readout, re-meeting the established faces at two of the four bodies, the electron and the muon, and parting at the other two, Theorem~229.1. These two chapters are where the Part is most honest about its reach: a periodic ring is an exact object of the saturation, but an atom with many electrons is a joined web whose admission the construction decides count by count, and the box lists the counts.

The reader should notice how the admission of a periodic web is decided, since the same two conditions return at every chapter of the Part. A ring is admitted when every join link lies within the valence bound of its docks and refused when a dock is already consumed; the parity of the cell count enters through the open dock that an odd ring leaves charged, Theorem~224.3, and the valence bound enters through the refusal of a second pairing at a consumed dock, Open result~226.2. The same two conditions decide the joined electrons at the probe in chapter 12-5, where the count of admitted electrons runs from one to ten and stops, Open result~228.1. Nothing in these decisions is a potential or a binding energy; they are statements about which pairs the grammar of the axiom allows at the docks of a copy, and the Part reads the whole of its subject out of them.

In the language of the literature, this Part names the lattice, the band structure, the dispersion relation, Cooper pairing, the Chern number, flux quantization, the periodic table and the many-body wavefunction. The correspondence is by name at each. The boxes establish periodic sinks, their readouts and the winding of a holonomy around a period.

The dependency lines of the Part reach back to the records, the probe and the saturation of Part I, to the energy chapter and the measurement of Part III, to the representations of Part V, to hydrogen and the lepton ladder of Part VII and to the stars of Part IX, whose docks are the sites at which the rings are joined. The lattice rests on the pairing move and the electron body, the band on the lattice and the step generator, the winding on the wash holonomy of Part XV, and the atoms on the probe and hydrogen. A reader following the derivation graph will see Part XIV as a Part that builds nothing new at the level of the link and instead joins the objects of the Parts ahead of it into rings and webs, which is why its theorems are the exact ones and its open results are the ones about admission.
\section{The lattice as the periodic sink}\label{mod:XII-1-lattice-as-the-periodic-sink}
\chdesig{14-1}{12-1}{XII-1-lattice-as-the-periodic-sink}{XII-1-lattice-as-the-periodic-sink}
\label{narr:open:12-1}The lattice is a periodic sink. Theorem~224.1 reads the chain of copies and the periodic readout, Theorem~224.2 the periodic network admitted by the grammar, Theorem~224.3 the periodic sink at a named cell count, three theorems, and Open result~224.4 states the periodic network of hub assemblies as an open result.
\begin{chtheorem}[The chain of copies and the periodic readout]\label{box:XII-1-lattice-as-the-periodic-sink:1}
A chain of two, three or four copies of the electron body joined at the dock by the pairing move, the join graph a path or a cycle, is one network whose direct record sum is an exact element at every one of the 5 arrays, every record graded five ways; the readout is periodic under the cycle shift at 2 of 2 cycles at four tests, so the lattice and the Bloch theorem are corollaries by name.
\end{chtheorem}
\chstanding{theorem}
\chdeps{the theorems of Chapters 1-0, 1-2, 1-4 to 1-6, 3-4, 7-10, 9-3, 9-6}{Theorem 225.1 and Theorem 226.1 and Open result 226.2}
\begin{chtheorem}[The periodic network admitted by the grammar]\label{box:XII-1-lattice-as-the-periodic-sink:2}
The periodic network is admitted by the grammar at both cells of the item, the library star with its dock (rings of N docked stars at three dock kinds, every join link a (2,2) pairing link between two docks, the translation an automorphism; 54 networks built, 17 periodic, 29 admitted) and the two-dock qutrit; the first link the dock pairing refuses is the like-kind link of the chain of three; the 17 periodic networks are absent from both truncated closures.
\end{chtheorem}
\chstanding{theorem}
\chdeps{the theorems of Chapters 1-0, 1-2, 1-4 to 1-6, 3-4, 7-10, 9-3, 9-6}{Theorem 225.1 and Theorem 226.1 and Open result 226.2}
\begin{chtheorem}[The periodic sink at a named cell count]\label{box:XII-1-lattice-as-the-periodic-sink:3}
The periodic sink exists at a named N with its cell an output: at the library star the sink of the ring of N docked stars under the pairing move at the docks is periodic at the even N two, four and six and refused at the odd N three and five (the charge carried at one open dock), its cell the lattice cell with its free dock paired at the half turn (sites 11, measure 91); at the two-dock qutrit the ring is its own sink at every N from two to six with the lattice limit cell exactly (sites 3, measure 16); 40 saturations terminated at 82 sinks, the parity law at every one.
\end{chtheorem}
\chstanding{theorem}
\chdeps{the theorems of Chapters 1-0, 1-2, 1-4 to 1-6, 3-4, 7-10, 9-3, 9-6}{Theorem 225.1 and Theorem 226.1 and Open result 226.2}
The three theorems are conditional on nothing beyond the text and the proved results they rest on. The third is the one derived in full below, the first two stated in their own boxes.
\begin{chwall}[The periodic network of hub assemblies]\label{box:XII-1-lattice-as-the-periodic-sink:4}
What stands open is the periodic network of hub assemblies, no configuration of the grammar at the join: the first link refused at N two is the (3,3) join link between two hub assemblies, neither end a dock qubit, while at the reading open beside the ring of three hub assemblies equals the tau network of the library by an explicit isomorphism, so the lepton of the third generation is the periodic network of three unit cells.
\end{chwall}
\chstanding{open result}
\chdeps{the theorems of Chapters 1-0, 1-2, 1-4 to 1-6, 3-4, 7-10, 9-3, 9-6}{Theorem 225.1 and Theorem 226.1 and Open result 226.2}
The odd N at the star is the object of the chapter of the chain of pairs and the dispersion the object of the next chapter.
\begin{chbridge}
The literature calls the link a maximally entangled pair of two qudits (exact), the dock a dangling edge at a terminal qubit (nearest), the join the composition of two tensor networks along shared sites (nearest), the sink a stable configuration (nearest), and the objects named the crystal, a periodic array of unit cells with the translation group as an exact symmetry \cite{AshcroftMermin1976,Kittel2005}, and the Bloch theorem \cite{Bloch1929}, by name only.
\end{chbridge}
\begin{chproof}[finite computation]
The premises are the Axiom and the six definitions of Section~\ref{mod:FRAMEWORK} (definitions), the link and its genres of Section~\ref{mod:I-2-link-and-its-six-genres}, the moves of Section~\ref{mod:I-4-moves-in-one-table}, the saturation of Section~\ref{mod:I-5-saturation-in-one-figure-seeds-waves-termination-sinks} and the read at a register of Section~\ref{mod:I-6-read-at-a-register-the-probe-and-the-measurement} (theorems), the wash at one act of Section~\ref{mod:II-4-measurement-the-wash-at-one-act-and-the-povm}, the lepton ladder of Section~\ref{mod:V-10-lepton-ladder-as-hub-assemblies}, the nucleon of Section~\ref{mod:VII-3-nucleon-structure-and-the-form-factor} and the hadron spectrum of Section~\ref{mod:VII-6-hadron-spectrum-and-the-mesons} (theorems). Nothing outside that list is used, and the chain from the axiom to this result is eight chapters deep.
The inputs are the two cells. The library star with its dock, 11 sites of dimensions $3,3,3,3,2,2,2,4,2,2,2$ at 11 links of measure 85 with the three docks $n_1$, $n_2$, $n_3$, and the two-dock qutrit, 3 sites of dimensions $2,2,3$ at 2 links of measure 12 with the docks $s_0$, $s_1$. The procedure is this. First, for a cell, a dock kind and a count $N$ from two to six, place $N$ copies of the cell and join copy $k$ to copy $k+1$ (and copy $N$ to copy one) by one $(2,2)$ pairing link between the two dock-labeled pendant qubits of the kind, so that the ring of $N$ docked cells has
\begin{equation}
\text{sites} = 11N, \qquad \text{links} = 11N + N = 12N \quad (\text{the star}); \qquad \text{sites} = 3N, \quad \text{links} = 3N \quad (\text{the qutrit}),
\end{equation}
and test its admission (no double link, every site at the valence bound, every join link between two docks with spare, the translation an automorphism by the explicit permutation of the sites). Second, saturate the admitted ring under the pairing move at the docks, wave by wave up to isomorphism, until no pairable dock pair remains (the stop condition), at two readings of the pairable set. Third, test every sink. Periodic if and only if the translation of the ring is an automorphism of the sink and the same cell stands at every position. Fourth, aggregate the letters per cell and $N$ and read the cell of every periodic sink. The totals at the checkpoints: 20 rings built and admitted (15 at the star at the three dock kinds, 5 at the qutrit), 40 saturations terminated at 82 sinks (58 dark, 24 charged), every sink at the stop condition and at the valence bound and every added link a $(2,2)$ link between two dock qubits. The translation carries 28 of the 82 sinks and the same cell stands at every position at 40, so 28 are periodic, 14 of the 41 sinks at the reading of reference. The smallest instance is the star at $N = 2$ at the kind $n_1$-$n_2$: 22 sites and 24 links with the two join links $c_1 n_1$-$c_2 n_2$ and $c_2 n_1$-$c_1 n_2$ and two open docks $c_1 n_3$, $c_2 n_3$. One wave adds the one pairing link $c_1 n_3$-$c_2 n_3$ and stops, the sink having 25 links and measure 182, dark, its translation an automorphism and the same cell at both positions, so periodic, the cell being the star with its free dock paired at the half turn. At $N = 3$ the ring has 33 sites, 36 links and three open docks, one wave pairs two of them and leaves $c_3 n_3$ open, the sink charged and refused. The parity law follows. The join leaves one free dock per copy, $N$ in all, and the pairing move pairs them in twos, so
\begin{equation}
\text{open docks at the sink} = N \bmod 2, \qquad \text{periodic} \iff N \equiv 0 \pmod 2 \ (\text{the star}),
\end{equation}
met at 40 of 40 saturations, while at the qutrit every dock is consumed by the join and the ring is its own sink at every $N$ from two to six, its cell the limit cell with 3 sites and measure 16. The result is the third theorem. The periodic sink exists at the even $N$ of the star and at every $N$ of the qutrit, the first named at the star at $N = 2$.
The check is by the counter-tests. The pairing move struck moves 15 of the 20 saturations, the first at the star at $N = 2$. The translation struck at crossed dock labels moves 19 of 20, the first at the qutrit at $N = 3$. The text struck leaves the gate unmet. Fourteen controls hold, the two cells equal to the unit-cell table and the twelve rings first built equal to the link table among them.
The result was reproduced by an independent route, two engines written apart agreeing at every ring, sink and letter.
\chcert{82}
\end{chproof}

\par\medskip\noindent{\itshape The lattice built, the next chapter reads the band at it.}\label{narr:close:12-1}\par
\section{The band and the dispersion}\label{mod:XII-2-band-and-the-dispersion}
\chdesig{14-2}{12-2}{XII-2-band-and-the-dispersion}{XII-2-band-and-the-dispersion}
\label{narr:open:12-2}Theorem~225.1 reads the band at the periodic web, a theorem resting on the records, the probe, the energy chapter of Part III and the lattice.
\begin{chtheorem}[The band at the periodic web]\label{box:XII-2-band-and-the-dispersion:1}
At the periodic networks of the preceding chapter, the rings of N docked stars at the three dock kinds and of N docked qutrits, re-built by the same join at N two to six, every link is a register of the step generator tiled by the translation, and the cycle-order spectrum of the ring is the spectrum of the cell repeated at every translation eigenvalue: 20 networks built, 20 periodic, 20 tiled, the spectrum equal to the tiling and translation-equal at every register at every one, with no parting N or register at either cell. The library star carries 11 registers and the two-dock qutrit 2, the join link the (2,2) register at every ring; the band at the star is the set of orders one to five at every N with the multiplicity per cell constant across N, at the qutrit the set of orders one, two and four likewise; the limit per cell equals the limit table of the lattice result; and the dispersion reading is flat at 10 of 10 tables, with no value and no material constant.
\end{chtheorem}
\chstanding{theorem}
\chdeps{Definitions 1.1 to 1.6, Theorems 18.1 to 18.4 and Open results 18.5 and 18.6, Theorems 28.1 and 28.2, Open results 28.3 and 28.4 and Theorem 28.5, Theorem 44.1 and Theorems 224.1 to 224.3 and Open result 224.4}{no theorem beyond it at this release}
The theorem is conditional on nothing beyond the text and the proved results it rests on. What stands open beside it is the dispersion of the literature as a dependence of the energy on a wave number, at no result (the dispersion reading flat at the limit of the count, and the phonon gloss of the lattice result naming the dispersion as the clause it lacks), and the band read at the two cells alone.
\begin{chbridge}
The literature calls the register the tensor factors of the common support (nearest), the step generator the Hamiltonian that generates the step (nearest), the cycle-order spectrum the spectrum of a cyclic shift (nearest), and the objects named the energy bands $E_n(k)$ over the Brillouin zone with the Kronig-Penney model their solved case \cite{Bloch1929,KronigPenney1931} and the tight-binding band $E(k) = E_0 - 2t\cos(ka)$ \cite{AshcroftMermin1976}, by name only.
\end{chbridge}
\begin{chproof}[census]
The premises are the Axiom and the six definitions of Section~\ref{mod:FRAMEWORK} (definitions), the records and the Born count of Section~\ref{mod:I-3-records-reads-and-the-born-count-the-gain-lemma}, the read at a register of Section~\ref{mod:I-6-read-at-a-register-the-probe-and-the-measurement} and the first-power operator and the step of Section~\ref{mod:II-6-energy-and-the-hamiltonian-the-first-power-operator-and-the-step} (theorems), and the lattice as the periodic sink of Section~\ref{mod:XII-1-lattice-as-the-periodic-sink} (a theorem). Nothing outside that list is used, and the chain from the axiom to this result is nine chapters deep.
The set counted is the 20 periodic networks, the rings of $N$ docked stars at the three dock kinds and of $N$ docked qutrits at $N$ from two to six, the 12 at $N$ two to four read from the lattice tables and re-built by the same join, the 8 at $N$ five and six built by the same join with the admission re-met. From a link of a ring, by its genre, we obtain a register of the step generator, the transport permutation the transpose of the register, and its cycle-order spectrum with its orders and its exponents at the common ring of one hundred and twenty. At $(2,3)$ the cycle type $1^{2} 4^{1}$ of order four, at $(2,2)$ the type $1^{2} 2^{1}$ of order two, at $(3,3)$ the type $1^{3} 2^{3}$, at $(2,4)$ the type $1^{2} 3^{2}$ of order three and at $(3,4)$ the type $1^{2} 5^{2}$ of order five,
\begin{equation}
(2,3): \ \text{orders } 1\!:\!2,\ 4\!:\!1, \quad \text{exponents } 0\!:\!3,\ 30\!:\!1,\ 60\!:\!1,\ 90\!:\!1 \ \text{ at the ring } 120 .
\end{equation}
From the translation of the ring, by its permutation of the sites, we obtain every register carried to a register of the same genre at the next copy at 20 of 20 networks, the first untiled register none. From the registers of a cell summed, by the tiling at every translation eigenvalue, we obtain the spectrum of the ring as $N$ copies of the cell spectrum with the join-link spectrum beside. At the star at $N = 2$ the orders
\begin{equation}
1\!:\!54, \ 2\!:\!22, \ 3\!:\!12, \ 4\!:\!6, \ 5\!:\!4 \quad = \quad 2 \cdot (1\!:\!25, \ 2\!:\!10, \ 3\!:\!6, \ 4\!:\!3, \ 5\!:\!2) \ + \ 2 \cdot (1\!:\!2, \ 2\!:\!1) ,
\end{equation}
the two copies and the two $(2,2)$ join links, and at $N$ three to six the multiplicities $81, 108, 135, 162$ at the order one scaling with $N$, equal to the tiling at every $N$ and translation-equal at every register. At the qutrit the orders one, two and four at the multiplicities $4, 2$ per cell and $2, 1$ per join link. From the cell with one pairing link, by the limit per cell, we obtain the limit pattern of the star, 11 sites, 12 links, measure 89 and register sum 28, with the spectrum $1\!:\!27, 2\!:\!11, 3\!:\!6, 4\!:\!3, 5\!:\!2$ at ten levels, and of the qutrit, 3 sites, 3 links, measure 16 and register sum 7, with the spectrum $1\!:\!6, 2\!:\!1, 4\!:\!2$, equal to the limit table of the lattice result at both. From the orders indexed by the translation eigenvalue at the $N$-th roots, by the dispersion reading, we obtain the same multiset at every root at 10 of 10 tables, flat, with the first parting root none. The property tested, the spectrum of the ring equal to the tiling of the cell spectrum and translation-equal at every register, holds at 20 of 20, which is the result.
The check is by the counter-tests. The translation struck at the first copy, a dock exchanged with a site of another dimension, moves the tiling at 20 of 20 networks. One dock of the second copy struck to dimension three moves 8 of 8 deciding cells and the face falls. The text struck leaves the gate unmet. Twelve controls hold, the twelve rings first built re-met at the lattice tables among them.
The result was reproduced by an independent route, two engines written apart agreeing at every register and spectrum.
\chcert{20}
\end{chproof}

\par\medskip\noindent{\itshape From single bodies at the lattice the thread passes to pairs.}\label{narr:close:12-2}\par
\section{The chain of pairs and the charged periodic sink}\label{mod:XII-3-chain-of-pairs-and-the-charged-periodic-sink}
\chdesig{14-3}{12-3}{XII-3-chain-of-pairs-and-the-charged-periodic-sink}{XII-3-chain-of-pairs-and-the-charged-periodic-sink}
\label{narr:open:12-3}Theorem~226.1 reads the chain of pairs at the periodic web, a theorem, and Open result~226.2 states the two open results, the pairing of two charged bodies and the charged periodic sink.
\begin{chtheorem}[The chain of pairs at the periodic web]\label{box:XII-3-chain-of-pairs-and-the-charged-periodic-sink:1}
At every ring of the lattice (the 12 rings of the lattice tables and 16 more built at N five to eight), the N join links are N (2,2) pairing links between dock-labelled pendant qubits of consecutive copies, disjoint and carried to the next by the translation, so the chain of N pairs is admitted at all 28 rings with no parting N or register; every pair register is 4 with the pair ket a fixed point of the step, the chain step is the product of the N pair steps at order two with the translation commuting with it, and the ring holonomy at the cycle is trivial at the letters Z and X everywhere and at the letter XZ trivial at every even N and nontrivial at every odd N; the photon chain precedent is re-met at the same pairing move.
\end{chtheorem}
\chstanding{theorem}
\chdeps{Definitions 1.1 to 1.6, Theorems 9.1 to 9.7 and Open result 9.8, Theorem 44.1, Theorem 150.1, Theorems 224.1 to 224.3 and Open result 224.4, Open results 241.1 and 241.2, Theorems 241.3 and 241.4 and Open results 241.5 to 241.7 and Open results 245.1 and 245.2, Theorem 245.3, Open results 245.4 to 245.8 and Theorem 245.9}{no theorem beyond it at this release}
The theorem is conditional on nothing beyond the text and the proved results it rests on, and the two open results below are stated beside it.
\begin{chwall}[The two open results: the pairing of two charged bodies and the charged periodic sink]\label{box:XII-3-chain-of-pairs-and-the-charged-periodic-sink:2}
What stands open is first the pairing of two charged bodies at the dock qubits: at the two-dock qutrit at N two the link between the electron dock and the first dock of the first cell is refused at the valence bound, every dock of every cell already consumed by the ring join, while at the library star two library electrons paired at the dock pairing are admitted at every N of the item (9 of 9 paired states, 72 of 72 lattice pairs at every distance to N seven), the paired state a sink of the pairing move and the two charged bodies paired away at 93 of 93 sinks of the 40 seeded saturations; and second the charged periodic sink, met at no cell: at the library star the charge is paired away at N two at every ring kind (one dark sink) and carried to the sink at N three (two charged sinks, each with one open dock end), the parity of the open docks deciding, while at the two-dock qutrit the attachment is refused at every N built; 16 attachments (12 admitted, 4 refused), 60 saturations terminated at 78 sinks (48 charged, 30 dark).
\end{chwall}
\chstanding{open result}
\chdeps{Definitions 1.1 to 1.6, Theorems 9.1 to 9.7 and Open result 9.8, Theorem 44.1, Theorem 150.1, Theorems 224.1 to 224.3 and Open result 224.4, Open results 241.1 and 241.2, Theorems 241.3 and 241.4 and Open results 241.5 to 241.7 and Open results 245.1 and 245.2, Theorem 245.3, Open results 245.4 to 245.8 and Theorem 245.9}{no theorem beyond it at this release}
At one face each open result would meet its object at a cell, the pairing admitted at both cells and a charged periodic sink at some cell. At the faces this chapter inherits the pairing is refused at the qutrit and the charge is paired away or carried to an open dock at the star, the parity deciding. The Chern number is the object of the next chapter.
\begin{chbridge}
The literature calls the pair a maximally entangled pair of qudits (exact), the pair ket a Bell state (exact), the holonomy the parallel transport around a closed path (nearest), the register the tensor factors of the common support (nearest), and the objects named the dimerized chain of the Peierls instability and the Su-Schrieffer-Heeger model \cite{Peierls1955,SuSchriefferHeeger1979} and the extra charge of a periodic chain, bound or carried as a band state \cite{AshcroftMermin1976}, by name only.
\end{chbridge}
\begin{chproof}[finite computation]
The premises are the Axiom and the six definitions of Section~\ref{mod:FRAMEWORK} (definitions), the link and its genres of Section~\ref{mod:I-2-link-and-its-six-genres}, the moves of Section~\ref{mod:I-4-moves-in-one-table}, the saturation of Section~\ref{mod:I-5-saturation-in-one-figure-seeds-waves-termination-sinks}, the first-power operator and the step of Section~\ref{mod:II-6-energy-and-the-hamiltonian-the-first-power-operator-and-the-step} and the photon of Section~\ref{mod:V-8-photon-docking-the-join-face-and-pair-production} (theorems), and the lattice of Section~\ref{mod:XII-1-lattice-as-the-periodic-sink} (a theorem). Nothing outside that list is used, and the chain from the axiom to this result is ten chapters deep.
The inputs are the two cells of the lattice (the star of measure 85 with the docks $n_1$, $n_2$, $n_3$, the two-dock qutrit of measure 12 with the docks $s_0$, $s_1$) and the rings of $N$ docked cells at $N$ from two to eight, 28 in all, the 12 at $N$ two to four re-met at the lattice tables column by column and the 16 at $N$ five to eight admitted by the same join. The procedure is this. First, locate the chain. The $N$ join links of the ring, each a $(2,2)$ pairing link between the dock-labeled pendant qubits of consecutive copies, disjoint, threaded by the cell path from one dock to the next (from $n_2$ through $Q_2$, $J$ and $Q_1$ to $n_1$ at the star, four links of the genres $(2,3)$, $(3,3)$, $(3,3)$, $(2,3)$), and test that the translation carries each pair to the next. Second, read the pair register of every pair as the product of its two dock dimensions and the step at the pair as the permutation of the gate-set row, and compose the chain step as the product of the $N$ pair steps. Third, compose the holonomy around the ring cycle on the three cap letters of the start dock at the ring of twenty-four. Fourth, aggregate per ring. At every pair of every ring
\begin{equation} \adjustbox{max width=\dimexpr\linewidth-3em\relax}{$\displaystyle \begin{gathered} \text{register} = 2 \cdot 2 = 4, \qquad \text{pair ket} = (1, 0, 0, 1), \\ \text{step: cycle type } 1^{2} 2^{1}, \ \text{order } 2, \ \text{eigenphases } 0\!:\!3,\ 12\!:\!1 , \end{gathered}$} \end{equation}
the two product states fixed and the two mixed states exchanged, the pair ket a fixed point. The smallest instance is the star at $N = 2$. Two pairs $P_1 = c_1 n_1$-$c_2 n_2$ and $P_2 = c_2 n_1$-$c_1 n_2$, the chain register $4^{2} = 16$, the chain step of cycle type $1^{4} 2^{6}$ and order two with the eigenphases $0\!:\!10, 12\!:\!6$, the generator at the three levels $0, 12, 24$ with the multiplicities $9, 6, 1$ at the spacing twelve, the translation commuting with the step and carrying $P_1$ to $P_2$. At $N = 3$ the chain register is 64 with the step of type $1^{8} 2^{28}$ and the generator at four levels $27, 27, 9, 1$, and at $N = 4$ the register 256 with the type $1^{16} 2^{120}$ and five levels $81, 108, 54, 12, 1$, the pattern
\begin{equation}
\text{register} = 4^{N}, \qquad \text{levels} = N + 1 \text{ at the spacing } 12, \qquad \text{multiplicities} = \binom{N}{j} 3^{N-j} .
\end{equation}
The holonomy around the cycle composed with the transport of the holonomy chapter is closed and a unitary on the register at every ring and letter. Trivial at the letters $Z$ and $X$ at all 28, and at the letter $XZ$ trivial at the 16 rings of even $N$ (the spectrum $z^{0}$ twice) and nontrivial at the 12 of odd $N$ (the spectrum $z^{0}$ and $z^{12}$, the returned set permuted setwise), the parity of the cycle length $5N$ deciding. The result is the chain at the pair register at 28 of 28 rings, the first parting none, and the two open results are stated in their own box.
The check is by the counter-tests. The translation struck at the first copy moves 244 cells, the chain first. The ring cut at one pair moves 1298 cells, the chain, the deciding cell and the holonomy among them. The text struck leaves the gate unmet. Twelve controls hold, the twelve rings first built re-met at the lattice tables among them.
The result was reproduced by an independent route, two engines written apart agreeing at every ring, register and holonomy.
\chcert{28}
\end{chproof}

\par\medskip\noindent{\itshape The next chapter reads the winding of the holonomy around a period.}\label{narr:close:12-3}\par
\section{The instrument tables of this Part}\label{mod:XIV-L-the-instrument-tables-at-the-repository}
\chdesig{14-L}{}{XIV-L-the-instrument-tables-at-the-repository}{XIV-L-the-instrument-tables-at-the-repository}
\begin{chtable}[The instrument tables of this Part]\label{box:XIV-L-the-instrument-tables-at-the-repository:1}
This Part lists 2 instrument chapters, printed among the instruments at the end of this compendium as chapters~\ref{mod:XII-4-chern-number} to~\ref{mod:XII-4-D538575-flxf-t1}, and the table beneath names each one with its key cells.
\end{chtable}
\chstanding{instrument}
\chdeps{Definitions 1.1 to 1.6}{no theorem beyond it at this release}
The listing moves no cell and states no result of its own. Each instrument stands at the repository as its chapter stood at the previous release of this compendium, and its tables are read there.
\begin{chbridge}
The tables listed here are instruments of the compendium, and this listing draws no correspondence with the literature.
\end{chbridge}
\begin{chproof}[instrument]
The listing in full. For each instrument chapter, printed among the instruments at the end of this compendium, the table gives its title, the result boxes it carries, the standing of each box with its citation, and the count of objects its certificate sentence names where it prints one.
{\small
\begin{longtable}{p{38mm}p{63mm}p{30mm}p{14mm}}
\hline
chapter & results & standing & objects \\
\hline
\endhead
The Chern number & The Chern number, awaited; The winding of the wash holonomy around one period; The winding around the full cycle of a joined periodic network; The flux quantization premise re-posed & instrument; theorem; theorem; open result & -- \\
The Chern number & The flux clause re-posed at the face & theorem & 45739 \\
\hline
\end{longtable}
}
The rule that produced it. Every chapter of this Part whose first standing beyond a placeholder reads instrument is printed among the instruments at the end of this compendium with its tables, and its key cells are the titles of its boxes, its standing lines and its certificate count, read from its text as it stood at the previous release.
No derivation stands in this chapter; the tables stand at the repository.
\emph{Check.} Every file listed was re-read at the repository and read the same as at the previous release, and the listing was re-cut from the files by a second route.
Nothing stands open at this listing. At this release two results of this Part are re-graded at their chapters: the periodic table (the chapter 14-5) re-posed under the full move set; atoms (the chapter 14-6), the probe read closed as a theorem at the second route.
\end{chproof}
\label{chend:XIV-L-the-instrument-tables-at-the-repository}
\section{The periodic table as the joined-electron sinks}\label{mod:XII-5-periodic-table-as-the-joined-electron-sinks}
\chdesig{14-5}{12-5}{XII-5-periodic-table-as-the-joined-electron-sinks}{XII-5-periodic-table-as-the-joined-electron-sinks}
\label{narr:open:12-5}Open result~228.1 states the periodic table as the joined-electron sinks, an open result resting on the moves, the saturation, the probe, the corollaries of Part III and the hydrogen chapter of Part VII.
\begin{chwall}[The periodic table as the joined-electron sinks]\label{box:XII-5-periodic-table-as-the-joined-electron-sinks:1}
At the reading taken as the one of reference (the probe at its three docks, the pairable set the docks, every arrangement, the saturation across copies, the class the component of the order parameter), the join of N library electrons at the probe is admitted at every N from one to ten and refused at none, every joined network reaches its sinks within the budget (425 sinks of 160 joined-network classes, the classes per N growing from 2 at one electron to 36 at ten and the sinks per N from 3 to 100), and the class of the sinks is charged at every N; the period read as the hub count of the sinks runs from 2 at one electron to 11 at ten, one hub per electron and one at the probe; the exclusion at the join is the valence law at a dock qubit, a second electron at a paired dock refused (0 refusals at one electron, 1080 at ten), read against the symmetrization postulate, a theorem; the bond is admitted at one plus one and at two plus two electrons, both dark; and the one-electron control re-meets the hydrogen cells of the hydrogen probe at 12 of 12. The class set does not change over N two to ten, so the periodic table by name stands at the joined sinks as a naming, an open result at that step, with no shell or row of the literature asserted; the alternative readings part at their first refusal or carry a class set beyond two, every one printed and none chosen.
Re-posed under the full move set: the joined-electron webs re-built at the re-cut catalogue electron re-meet the classes and the chains above at ten of ten counts of electrons and the sinks and the rows at none of ten, the join admitted at every count with the classes per count 2, 4, 6, 9, 12, 16, 20, 25, 30 and 36 and the sinks per count 3, 8, 14, 22, 31, 42, 54, 68, 83 and 100; under the twenty-eight moves beside the pairing move alone every closure is unstopped at the candidate bound and the sink of reference is among the closure webs at no count, the first row parting at the cell type of its first sink, neutral against charged, every row parting; none of the twenty-eight moves pairs two docks at the joined web at one electron, so the dock pairing of the reference reading is a move of the table at no application of 106, every dock end open where the join of the table asks closed and equal ends.
\end{chwall}
\chstanding{open result}
\chdeps{Definitions 1.1 to 1.6, Theorems 28.1 and 28.2, Open results 28.3 and 28.4 and Theorem 28.5, Theorem 47.1, Theorem 143.1, Open results 241.1 and 241.2, Theorems 241.3 and 241.4 and Open results 241.5 to 241.7 and Open results 245.1 and 245.2, Theorem 245.3, Open results 245.4 to 245.8 and Theorem 245.9}{no theorem beyond it at this release}
At one face the joined-electron sinks would exist at every N of the budget with a class structure that changes at definite N. At the face this chapter inherits the class set is charged at every N from one to ten and the change set over two to ten is empty, the first named the class never changing. The chains of electrons at the library star are the attachments of the charged periodic sink and the pairs of charged bodies of the preceding chapter, read there.
\begin{chbridge}
The literature calls the dock a dangling edge at a terminal qubit (nearest), the sink a stable configuration (nearest), the probe the measuring body at the readout (nearest), the valence the degree of a vertex (nearest), and the objects named the periodic table of Mendeleev with the exclusion principle of Pauli and the shell closures \cite{Mendeleev1869,Pauli1925} and the atom of Bohr \cite{Bohr1913}, by name only. The construction supplies the valence law at a dock qubit at Section~\ref{mod:I-2-link-and-its-six-genres}, at which the symmetrization postulate is located.
\end{chbridge}
\begin{chproof}[census]
The premises are the Axiom and the six definitions of Section~\ref{mod:FRAMEWORK} (definitions), the moves of Section~\ref{mod:I-4-moves-in-one-table}, the saturation of Section~\ref{mod:I-5-saturation-in-one-figure-seeds-waves-termination-sinks} and the read at a register of Section~\ref{mod:I-6-read-at-a-register-the-probe-and-the-measurement} (theorems), the quantum-mechanical corollaries of Section~\ref{mod:II-9-qm-corollaries-the-fifteen-entries} (a theorem) and the hydrogen probe of Section~\ref{mod:V-4-hydrogen-the-rydberg-ladder-and-the-fine-structure} (a theorem). Nothing outside that list is used, and the chain from the axiom to this result is nine chapters deep.
The set counted is the joined networks of $N$ library electrons at the probe at $N$ from one to ten at the reading of reference (the probe at its three docks, the pairable set the docks, every arrangement, the saturation across copies), 160 joined-network classes in all. The electron of the library has 9 sites and 10 links of measure 79 with the docks $w_1$, $w_2$ and the first power 67 at the dock spelling. The probe has three docks. From $N$ electrons and the probe, by the dock pairing at every pairing the pairable sets admit, we obtain the joined networks of the count $N$, admitted at every $N$ from one to ten and refused at none. From each, by the pairing move at the stop condition (no pairable pair), we obtain its sinks, every one at the stop condition and the valence bound, the budget reached at no network, with the counts
\begin{equation}
\text{sinks per } N: \ 3, 8, 14, 22, 31, 42, 54, 68, 83, 100 \quad (N = 1, \ldots, 10), \qquad 425 \text{ in all} ,
\end{equation}
and the classes per $N$ growing from 2 to 36. From a sink, by its open dock ends, we obtain its class, and the property tested, a change of the class set at some $N$ from two to ten, fails. The class set is charged at every $N$ from one to ten, the change set empty, so the face inherited is the one stated and the first parting element is the class never changing at $N$ one to ten. From a sink, by its hub count, we obtain the period, running from 2 at one electron to 11 at ten, one hub per electron and one at the probe. From the join of a second electron at a paired dock, by the valence law at a dock qubit, we obtain the exclusion, 0 refusals at one electron rising to 1080 at ten. The smallest instance is one electron at the probe. Two joined networks (the electron at the first or at the third dock), three sinks in all, each adding one link from a probe dock to the electron dock $w_2$ and leaving one probe dock open, charged, with the order parameter
\begin{equation}
(\text{class}, W) = (\text{charged}, 214990848), \qquad \text{hubs } 2, \quad \text{first power at the dock spelling } 166 = 67 + 99 ,
\end{equation}
the hydrogen cells of the hydrogen probe re-met at 12 of 12. The bond of one electron plus one is admitted and dark, with the first power 348 at the dock spelling. The readings beside part at their first refusal (every electron at the probe alone refused at $N = 4$ for want of spare) or carry the same charged class set, none chosen.
The check is by the counter-tests. The pairing move struck moves the sinks at 10 of 10 counts, the first at one electron. The pairable set of the probe struck moves the joins at 10 of 10. The text struck leaves the gate unmet. Fourteen controls hold, the electron re-met at the seeds table and at the fixed library among them.
The result was reproduced by an independent route, two engines written apart agreeing at every joined network, sink and class.
\chcert{425}
\end{chproof}

\par\medskip\noindent{\itshape The Part closes on the probe web of hydrogen and its many-body readout.}\label{narr:close:12-5}\par
\section{Atoms: the hydrogen probe web and the N-body readout}\label{mod:XII-6-atoms-the-hydrogen-probe-web-and-the-n-body-readout}
\chdesig{14-6}{12-6}{XII-6-atoms-the-hydrogen-probe-web-and-the-n-body-readout}{XII-6-atoms-the-hydrogen-probe-web-and-the-n-body-readout}
\label{narr:open:12-6}Theorem~229.1 states the probe read at the four bodies joined to the measuring body as an open result, resting on the records, the probe, the measurement of Part III, the representations of Part V, the hydrogen chapter and the lepton ladder.
\begin{chtheorem}[The probe read at the four bodies joined to the measuring body]\label{box:XII-6-atoms-the-hydrogen-probe-web-and-the-n-body-readout:1}
The four bodies joined to their measuring bodies at the attachment site, with the joined read opened at the dock or nowhere under the transport of reference, give the unit and the register as outputs of the readout, and every one of those readings is equal at two routes: the unit and the register at the three unit readings of the eleven joined webs and of the six body rows at the two ladders, fifty-one readings in all, re-derived at a second route from the re-cut catalogue of webs, the nucleon figure, the derived-ladder readings and the transport table, read equal at the two routes at 51 of 51, 33 of 33 at the web level and 18 of 18 at the body level. At 2 of 4 bodies the outputs re-meet the established faces, the electron and the muon at the dock at the register 16 against the dock-struck readings 67 and 107 at 16, and at 2 of 4 they part, the tau first, its joined read opening at neither the hub nor the dock, both pendant qubits of the library tau hanging on (2,2) links to hub leaves closed at both ends, and the proton beside, the join attaching the qubit end of the photon word at the core, a qutrit site that is the far end of no link read from its larger end; the hub unit is reached by no joined read, so the hub-or-dock question of the previous releases resolves to the dock wherever the joined read opens at all. The first parting of the previous open result, the tau at the library ladder, re-reads at the second route as the same output, neither at no register against the dock-struck reading 147 at the register 16, a parting of the inherited face and not of the construction, which is thereby closed; at the derived ladder the tau carries the two (2,3) docks of the electron and re-meets the dock at 16, the ladder reading 3 of 4; the deciding table is re-derived in Lean.
\end{chtheorem}
\chstanding{theorem}
\chdeps{Definitions 1.1 to 1.6, Theorems 18.1 to 18.4 and Open results 18.5 and 18.6, Theorems 28.1 and 28.2, Open results 28.3 and 28.4 and Theorem 28.5, Theorem 42.1, Open results 66.1 to 66.3, Theorem 66.4 and Open result 66.5, Theorem 143.1 and Theorem 152.1}{no theorem beyond it at this release}
At one face the unit and the register would be outputs at four of four bodies re-meeting the faces the office results carried. At the face this chapter inherits they re-meet at two and part at the tau first, the parting named with both objects below. The hydrogen atom of this compendium is the hydrogen probe network at Part VII, the mass ratio not an input. The N-body readout is the readout of the many-body join at which the joined electrons of the preceding chapter are read.
\begin{chbridge}
The literature calls the far end the partner qudit of an entangled pair (exact), the open end a dangling edge (exact), the dock a dangling edge at a terminal qubit (nearest), the readout the outcome probabilities of a measurement (exact), and the objects named the atom of Bohr and of the wave equation \cite{Bohr1913,Schrodinger1926} and the many-electron atom of the self-consistent field \cite{Hartree1928,Fock1930}, by name only. The construction supplies the first-power operator at Section~\ref{mod:II-6-energy-and-the-hamiltonian-the-first-power-operator-and-the-step}.
\end{chbridge}
\begin{chproof}[census and exact cell]
The premises are the Axiom and the six definitions of Section~\ref{mod:FRAMEWORK} (definitions), the records and the Born count of Section~\ref{mod:I-3-records-reads-and-the-born-count-the-gain-lemma}, the read at a register of Section~\ref{mod:I-6-read-at-a-register-the-probe-and-the-measurement}, the wash at one act of Section~\ref{mod:II-4-measurement-the-wash-at-one-act-and-the-povm}, the representations from the open-end census of Section~\ref{mod:IV-3-representations-from-the-open-end-census} (an open result), the hydrogen probe of Section~\ref{mod:V-4-hydrogen-the-rydberg-ladder-and-the-fine-structure} and the lepton ladder of Section~\ref{mod:V-10-lepton-ladder-as-hub-assemblies} (theorems). Nothing outside that list is used, and the chain from the axiom to this result is nine chapters deep.
\emph{The exact part.} The object is the electron body joined to its measuring body, the photon word, at the attachment site, the dock $w_1$: 26 sites and 28 links, two hubs, the docks $w_1$ and $w_2$, measure 206. From the joined network, by the direct two-exchange sum (the ordered pairs of one-exchange acts at the unstruck links, each act at the order one the chain-sum register weight multiplied by the dressing), we obtain the readout through the joined network as an exact element in the basis one, root two, root three,
\begin{equation}
R_{\mathrm{whole}} = \frac{8\ 990\ 453\ 214\ 530\ 839}{2\ 949\ 120\ 000} + \frac{661}{225}\sqrt{2} + \frac{5977541}{147456000}\sqrt{3} ,
\end{equation}
the pair sum equal to the first power squared, and the first power 1746 of the joined network with the probe docks struck as at the hydrogen network. From every link end of the joined network, by the transport of reference (the far end open if and only if the transported far set is a frame, at the ordered dimension pairs three to two, four to two and four to three and at no other), we obtain the open ends of the joined read, and at the attachment site $w_1$ the two links $L_{9}$ and $L_{27}$ are open, so the unit is the dock and the register is
\begin{equation}
16 = 2 \cdot 2 \cdot 2 \cdot 2 ,
\end{equation}
equal to the dock-struck reading 67 at the register 16 of the body read alone. At the muon the same procedure gives the dock at the register 16 against the dock-struck 107.
\emph{The census part.} The set counted is the four bodies of the library ladder with their measuring bodies of reference (six measurement networks, 11 joined networks built of 12). From each, by the open ends at the attachment site, we obtain the unit output and the register, and the property tested, the output re-meeting the face the office result carried, holds at 2 of 4, the electron and the muon at the dock at 16. The first parting element is the tau, whose joined read opens at neither the hub nor the dock (both pendant qubits of the library tau hang on $(2,2)$ links to hub leaves closed at both ends), the unit output neither at the register none against the face inherited, the dock-struck reading 147 at the register 16. The proton parts beside, the join attaching the qubit end of the photon word at the core, a qutrit site that is the far end of no link read from its larger end. At the derived ladder, where the tau carries the two $(2,3)$ docks of the electron, the tau re-meets the dock at 16 and the ladder reads 3 of 4, the proton the one parting. The two readings beside (the unit at the apparatus link alone, the unit at the open ends reached through the body) meet 1 of 4 each, printed and not chosen.
The check is by the counter-tests. The measuring body struck, the body read alone, moves 37 of 254 compared cells and the readout at 11 of 11 networks. An open end washed by hand moves 65 cells and the unit at 50 of 168 cells. The body weight placed otherwise moves the readout alone at 11 of 11. Fourteen controls hold at the fields of the networks, the sums and offices of the four bodies read alone re-met among them, one control re-typed at two presentation columns of a parent table at review.
The result was reproduced by an independent route, two engines written apart agreeing at every readout, unit and register.
\chcert{11}
\end{chproof}
\begin{figure}[h]
\begin{center}
\setlength{\unitlength}{1pt}
\begin{picture}(300,80)(0,0)
\put(60,40){\circle{26}}
\put(52,36){$e$}
\put(73,40){\line(1,0){40}}
\put(113,40){\circle*{6}}
\put(116,46){$w_1$}
\put(113,40){\line(1,0){60}}
\put(173,40){\circle*{6}}
\put(176,46){$g_6$}
\put(173,40){\line(1,0){40}}
\put(213,40){\circle{26}}
\put(205,36){$\gamma$}
\put(226,40){\line(1,0){40}}
\put(266,40){\circle*{6}}
\put(269,46){$Q_1$}
\put(47,40){\line(-1,0){25}}
\put(22,40){\circle*{6}}
\put(6,46){$w_2$}
\put(40,6){the electron body}
\put(190,6){the photon word}
\put(100,66){the attachment site}
\end{picture}
\end{center}
\caption{The electron body joined to its measuring body, the photon word, at the attachment site, the dock $w_1$. The joined read opens at the dock and the register at the readout is sixteen.}
\end{figure}

\par\medskip\noindent{\itshape Matter read at the lattice and the atom, Part XV turns to information and computation.}\label{narr:close:12-6}\par
